\subsection{Phát biểu bài toán}
Mạng ở thời điểm $t$ là đồ thị $G_t=(V,E_t)$ với $V$ là tập UAV. Mỗi cạnh
$e=(a,b)$ có chất lượng đo được $d_e(t)\in[0,1]$, là tỷ lệ nhận probe trong
cửa sổ một giây. Bộ ước lượng chất lượng tạo ra một đại lượng $x_e(t)$ dùng cho
quyết định, theo một trong hai họ:
\begin{align}
\text{đo và làm mượt:}\quad & x_e(t)=\alpha\, d_e(t)+(1-\alpha)\,x_e(t-1),\label{eq:ewma}\\
\text{dự báo học được:}\quad & x_e(t)=\hat g_e(t)\approx
\frac{1}{H}\sum_{\tau=t+1}^{t+H} d_e(\tau),\qquad H=5\ \text{giây}.\label{eq:fore}
\end{align}
Công thức~\eqref{eq:ewma} là EWMA một tham số; $\alpha=0{,}5$ là cấu hình của
các nghiên cứu trước và được dùng làm baseline của v8. Công thức~\eqref{eq:fore}
là bộ hồi quy GNN, học có giám sát từ dữ liệu lịch sử và chỉ dùng quan sát
nhân quả tới thời điểm $t$.

Hai hành động dùng chung $x$. Hành động định tuyến chọn đường
\begin{equation}
\pi_t=\arg\min_{P\in\mathcal{P}(s,d)}\ \sum_{e\in P}\frac{1}{\max(x_e(t),\varepsilon)},
\label{eq:etx}
\end{equation}
tức chi phí kiểu ETX~\cite{decouto} tính trên $x$ thay vì trên số đo thô, với
luật ghim tuyến có margin và dwell mô tả ở Mục~\ref{sub:v8}. Hành động bố trí
dùng chất lượng đường tốt nhất
\begin{equation}
Q(t)=\max_{P\in\mathcal{P}(s,d)}\ \min_{e\in P} x_e(t),
\label{eq:bottleneck}
\end{equation}
và kích hoạt một luật di chuyển relay cố định khi $Q(t)$ xuống dưới ngưỡng.

Phân biệt trung tâm của bài là giữa hai thước đo trên cùng một $x$. Chất lượng
\emph{dự báo} đo bằng sai số bình phương trung bình
$\mathrm{MSE}=\mathbb{E}\big[(x_e(t)-\frac{1}{H}\sum_\tau d_e(\tau))^2\big]$,
tính ngoại tuyến trên trace. Chất lượng \emph{dịch vụ} đo bằng PDR và số giây
chia cắt, tính trong vòng kín khi $x$ thật sự điều khiển~\eqref{eq:etx}
hoặc~\eqref{eq:bottleneck}. Câu hỏi của bài là hai thước đo này có xếp hạng các
bộ ước lượng theo cùng một thứ tự hay không.

\subsection{Môi trường và thiết kế thí nghiệm}
Thực nghiệm dùng ns-3.45~\cite{ns3}, Wi-Fi 802.11g và OLSR~\cite{rfc3626} trên chuỗi 5 hoặc 7 UAV.
Nguồn và đích ở hai đầu chuỗi; các nút còn lại làm relay. Mỗi episode dài
240 giây, luồng chính UDP 256 kbps, hai luồng nền 128 kbps. Hai bối cảnh gồm
static và gradual: đầu mút đứng yên hoặc di chuyển tách xa từ giây 30.
Binary và runtime kế thừa nghiên cứu trước, gồm baseline OLSR-ETX do nhóm
cài đặt và thời hạn ARP DeadTimeout 1 giây. Mọi arm dùng cùng binary,
kênh, lưu lượng và seed của khối tương ứng.
Vùng mô phỏng $1800\times1200$ m, độ cao 50--150 m. Khoảng cách ban đầu
giữa các nút liên tiếp danh định 24 m với jitter tối đa 0,72 m mỗi tọa độ.
Kênh dùng log-distance (số mũ 2,2, suy hao tham chiếu 40 dB tại 1 m) và
Nakagami ($m=3;1{,}5;1$ theo các dải khoảng cách dưới 96 m, 96--180 m,
trên 180 m). Công suất phát 20 dBm, tốc độ PHY cố định 6 Mbps. Tầm phủ
danh định 240 m dùng để đặt độ nhạy thu, không phải ngưỡng cắt liên kết
theo khoảng cách. Warmup 10 giây; probe broadcast cách 0,05 giây, tổng hợp
cửa sổ một giây; luồng chính gói 512 B, 15.000 gói mỗi episode.

Protocol được viết trước các mô phỏng mới. Mỗi pha đóng gói mã nguồn, mô
hình, seed plan và định danh runtime bằng SHA-256; test kế thừa đúng nguồn
và mô hình validation. Mỗi cell chạy một lần; queue bốn worker giữ lock
toàn server. Validation có 80 cell (hai static và tám gradual mỗi cấu hình,
bốn arm); test có 192 cell (bốn static và 20 gradual mỗi cấu hình, bốn arm).
Không chọn lại mô hình theo kết quả dịch vụ của validation hoặc test.

\subsection{Dự báo liên kết và đối chứng}
Mô hình dự báo tỷ lệ nhận probe trung bình trong năm giây tiếp theo, không
phải xác suất sống sót của liên kết. Dữ liệu phát triển gồm 48 trace HOLD
test v7 để fit và 20 trace HOLD validation v7 để chọn epoch. Các dữ liệu đã
quan sát này được coi hoàn toàn là dữ liệu phát triển cho nghiên cứu mới;
v8/v9 dùng seed độc lập đã kiểm toán. Mỗi episode có 236 gốc dự báo hợp lệ,
tạo 11.328 mẫu fit và 4.720 mẫu chọn epoch; các frame tương quan không được
coi là các lần lặp độc lập khi kiểm định hiệu quả điều khiển.

Đầu vào chỉ gồm quan sát hiện tại và quá khứ. Tám đặc trưng nút gồm tọa độ
tương đối với tâm đội hình, vận tốc, bậc hàng xóm và tỷ lệ nhận trung bình.
Bảy đặc trưng cạnh gồm khoảng cách, RSSI, tỷ lệ nhận hiện tại, EWMA hệ số
0,5, chênh tỷ lệ nhận, vận tốc tách xa và độ lớn vận tốc tương đối. Tọa độ
và khoảng cách chia 240 m; vận tốc chia 10 m/s; RSSI biến đổi thành
$(RSSI+90)/40$, giá trị thiếu mã hóa $-100$ dBm. Cạnh truyền thông điệp có
tỷ lệ nhận hiện tại ít nhất 0,1. Không đọc quỹ đạo nhiệm vụ trong tương lai.

Edge-SAGE gồm hai lớp tổng hợp trung bình với đặc trưng cạnh, 64 chiều và
ReLU; bộ giải mã cặp nút 64 chiều và sigmoid, đối xứng hóa logits. GNN có
19.137 tham số. MLP bỏ hai nhánh truyền thông điệp học được, còn 13.505 tham
số; vẫn giữ đặc trưng bậc nút. Vì thế ablation chưa tách riêng tác động của
truyền thông điệp khỏi khác biệt dung lượng mô hình. Cả hai dùng MSE, Adam
lr 0,001, weight decay $10^{-5}$, batch 128, clip norm 5, tối đa 60 epoch,
patience 10, seed 20261123. GNN/MLP chọn epoch 58/59. Sai số lớn nhất giữa
suy luận NumPy và PyTorch trên 40 snapshot kiểm tra là $1{,}19\times10^{-7}$.
Đây là học có giám sát; các nghiên cứu không dùng học tăng cường.

\subsection{Nghiên cứu định tuyến v8}\label{sub:v8}
Bốn arm là chất lượng hiện tại (\code{lq-route}), dự báo GNN
(\code{gnn-route}), dự báo MLP (\code{mlp-route}) và OLSR-ETX-fast.
Ba arm ghim tuyến giữ cùng luật ETX với margin 1,0, warmup ba bước và dwell
năm giây; kiểm tra liên kết hỏng vẫn dựa trên tỷ lệ nhận hiện tại dưới 0,1.
ETX-fast dùng HELLO 0,5 s, TC 1,25 s đã chọn ở v7. Relay đứng yên; tốc độ
mỗi đầu mút gradual có giá trị danh định 1 m/s, nhân hệ số [0,85; 1,15] theo
layout seed. Nguồn đi theo chiều $+y$, đích theo chiều $-y$; cùng seed giữ
cùng chuyển động đầu mút giữa các arm. Validation chỉ kiểm tra kỹ thuật: đủ cell,
đúng tuyến, cờ simulator, forecast, bộ đếm và luồng nền, không blackout khởi
động và không di chuyển relay. Go không phụ thuộc chênh PDR.

Ba bước kiểm định theo thứ tự cố định là PDR GNN trừ hiện tại, GNN trừ MLP,
rồi GNN trừ ETX-fast. Một bước đạt khi cận dưới CI 95\% dương; chỉ diễn giải
khẳng định khi mọi bước trước đạt. MSE test được đo riêng cho GNN, MLP và
EWMA trên cùng các trace \code{lq-route}, tách khỏi PDR vòng kín.

\subsection{Nghiên cứu vị trí v9}
V9 giữ nguyên trọng số v8 và giữ định tuyến theo chất lượng hiện tại ở mọi
arm. Các arm là relay đứng yên, kích hoạt di chuyển phản ứng, kích hoạt bằng
GNN và kích hoạt bằng MLP. Chất lượng đường tốt nhất là $Q=\max_P\min_{e\in P}q_e$.
Sau ba frame, reactive kích hoạt nếu EWMA hiện tại có $Q<0{,}70$ trong hai
frame liên tiếp; GNN/MLP dùng thêm điều kiện tương tự trên dự báo. Cạnh dự
báo chỉ được dùng nếu số đo hiện tại ít nhất 0,1.

Các relay được xếp theo hình chiếu lên đường nối đầu mút hiện tại rồi hướng
tới các điểm chia đều đoạn đó. Giới hạn: 8 m/s, 400 m mỗi relay, deadband
1 m, tối đa mười bước mỗi đợt và nghỉ năm bước. Dừng sớm sau ba frame có
chất lượng hiện tại và dự báo (nếu có) ít nhất 0,80. Kiểm tra khoảng cách
nhỏ nhất trong bước phải đạt 5 m; nếu không, lệnh chuyển động relay bằng 0.
Luật này giống nhau ở mọi arm có chuyển động, chỉ khác trigger.

Tốc độ danh định mỗi đầu mút gradual tăng lên 2 m/s, vẫn giữ hệ số ngẫu nhiên
theo seed như v8, để tạo thách thức chia cắt; khác phân phối fit 1 m/s.
Validation đòi kỹ thuật/an toàn đạt và ít nhất 8/16
baseline gradual đứng yên mất kết nối ít nhất năm giây. Nếu không đạt thì
dừng test, không chỉnh lại tốc độ trong cùng nghiên cứu.

Endpoint chính là số giây mất kết nối của reactive trừ GNN; sau đó xét PDR
GNN trừ reactive, rồi số giây mất kết nối của MLP trừ GNN. Mất kết nối là
nguồn và đích khác thành phần liên thông trong đồ thị probe ngưỡng 0,5, đếm
trên frame 1--240. Số giây OLSR không có tuyến và PDR là các đại lượng khác.
Đây là chia cắt theo ngưỡng đồ thị đo được, không phải chứng minh hoàn toàn
không có gói vô tuyến nào có thể truyền giữa các thành phần.
Quãng đường, outage dài nhất, số đợt di chuyển và PDR luồng nền là mô tả.

\subsection{Phân tích thống kê và chi phí}
Mỗi so sánh giữ bốn arm cùng khối seed; bootstrap 10.000 lần trên 40 khối
gradual test, 20 mỗi cấu hình, cân bằng n5/n7, RNG 20261120. Bootstrap~\cite{efron1993}
được áp dụng ở cấp kịch bản để giữ tương quan giữa các arm. Kết quả âm và
các bước bị đóng cổng vẫn được trình bày. CI chứa 0 không chứng minh tương
đương. Kết quả static chỉ mô tả với tám khối.

Collector đo thời gian toàn bộ quyết định trên máy ảo KVM tám vCPU (model
Intel Xeon Silver 4110 2,10 GHz), chạy bốn worker, mỗi worker giới hạn thư
viện tính toán về một luồng. Probe phát trong mọi arm; telemetry và lệnh tập trung chưa được
truyền thành gói cạnh tranh kênh. Không quy thời gian tính thành độ trễ điều
khiển đầu cuối, số byte thành airtime, hoặc quãng đường thành năng lượng.

Chi phí báo hiệu được đo ở nghiên cứu v7 với cùng binary và runtime, nên các
con số sau áp dụng cho cấu hình của v8: OLSR thuần khoảng 0,29 kB/s, OLSR-ETX
0,44 kB/s, OLSR-ETX với HELLO 0,5 giây 1,78 kB/s, và ghim tuyến khoảng
5,55 kB/s gồm OLSR dự phòng 0,29 kB/s, probe 5,02 kB/s và báo cáo về bộ điều
khiển ước lượng 0,25 kB/s. Đây là byte tầng IP cộng phần báo cáo ước lượng,
không quy đổi thành phần trăm dung lượng kênh. Probe được phát ở mọi arm nên
airtime của probe là chung giữa các arm, nhưng chỉ được hạch toán vào chi phí
của các phương pháp sử dụng probe. Phần báo cáo tập trung chưa chiếm kênh trong
mô phỏng. Vì vậy mọi so sánh giữa ghim tuyến và OLSR-ETX trong bài phải được
đọc như giới hạn trên của một bộ điều khiển tập trung lý tưởng hóa một phần.
